\documentclass[12pt,a4paper]{article}

% Encoding and fonts
\usepackage[utf8]{inputenc}
\usepackage[T1]{fontenc}
\usepackage{lmodern}

% Page geometry
\usepackage[margin=2.5cm, top=3cm, bottom=3cm]{geometry}

% Math
\usepackage{amsmath,amssymb,amsthm}

% Tables
\usepackage{booktabs}
\usepackage{array}

% Lists
\usepackage{enumitem}

% Headers
\usepackage{fancyhdr}

% Colors and links
\usepackage{xcolor}
\usepackage{hyperref}

% Verbatim
\usepackage{fancyvrb}

% Turkish support
\usepackage[turkish]{babel}

% Header/footer setup
\pagestyle{fancy}
\fancyhf{}
\fancyhead[L]{\textit{Optimizasyona Giri\c{s} -- Ders Notlar{\i}}}
\fancyhead[R]{\textit{Lecture 1}}
\fancyfoot[C]{\thepage}
\renewcommand{\headrulewidth}{0.4pt}
\renewcommand{\footrulewidth}{0.4pt}

% Theorem environments
\theoremstyle{definition}
\newtheorem{ornek}{\"Ornek}[section]
\newtheorem{tanim}{Tan{\i}m}[section]

\theoremstyle{remark}
\newtheorem*{yorum}{Yorum}

% Title
\title{
    \vspace{-1cm}
    {\LARGE \textbf{Optimizasyona Giri\c{s} -- Ders Notlar{\i}}} \\[0.5cm]
    {\Large Lecture 1: Optimizasyonun Temelleri} \\[0.3cm]
    {\large Giri\c{s}, Konveks Kavramlar, Kalk\"ul\"us ve Lineer Cebir Tekrar{\i}}
}
\author{}
\date{}

\begin{document}

\maketitle
\thispagestyle{fancy}

\tableofcontents
\newpage

%% ============================================================
%% SECTION 1
%% ============================================================
\section{Optimizasyon Nedir?}

\subsection{En Basit Tan{\i}m{\i}yla}

Optimizasyon = \textbf{en iyi de\u{g}eri bulmak}.

Bir fonksiyonun \textbf{minimum} (en k\"u\c{c}\"uk) veya \textbf{maximum} (en b\"uy\"uk) de\u{g}erini veren $x$'i ar{\i}yoruz.

\[
\min_{x} f(x)
\]

\textbf{\"Ornek:} $f(x) = x^2 - 4x + 5$ fonksiyonunun minimumunu bul.

Bildi\u{g}in t\"urevi al{\i}p s{\i}f{\i}ra e\c{s}itle:
\begin{itemize}
    \item $f'(x) = 2x - 4 = 0 \;\Rightarrow\; x = 2$
    \item $f(2) = 4 - 8 + 5 = 1$
\end{itemize}

$x = 2$'de minimum de\u{g}er $1$'dir. \.{I}\c{s}te optimizasyon bu kadar basit ba\c{s}l{\i}yor!

\subsection{Optimizasyon Probleminin Genel Formu}

\[
\min_{x} f(x) \quad \text{\"oyle ki (subject to)} \quad x \in \Omega
\]

\begin{itemize}
    \item \textbf{$f(x)$:} Ama\c{c} fonksiyonu (objective function) -- minimize etmek istedi\u{g}imiz \c{s}ey
    \item \textbf{$x$:} Karar de\u{g}i\c{s}keni (decision variable) -- de\u{g}i\c{s}tirebildi\u{g}imiz de\u{g}er
    \item \textbf{$\Omega$:} Uygun b\"olge (feasible region) -- $x$'in alabilece\u{g}i de\u{g}erler k\"umesi
\end{itemize}

\subsection{S{\i}n{\i}fland{\i}rma}

\subsubsection{Minimizasyon vs Maksimizasyon}
\begin{itemize}
    \item $\min f(x) = -\max(-f(x))$
    \item Yani birini \c{c}\"ozebildin mi, di\u{g}erini de \c{c}\"ozersin. Hep \textbf{minimizasyon} \"uzerinden gideriz.
\end{itemize}

\subsubsection{S\"urekli (Continuous) vs Ayr{\i}k (Discrete)}
\begin{itemize}
    \item \textbf{S\"urekli:} $x$ herhangi bir reel say{\i} olabilir ($x = 3.14159\ldots$)
    \item \textbf{Ayr{\i}k:} $x$ sadece tam say{\i} olabilir ($x = 1, 2, 3, \ldots$)
    \item \"Ornek: ``Ka\c{c} kg \c{s}eker alay{\i}m?'' $\to$ s\"urekli, ``Ka\c{c} araba \"ureteyim?'' $\to$ ayr{\i}k
\end{itemize}

\subsubsection{K{\i}s{\i}tl{\i} (Constrained) vs K{\i}s{\i}ts{\i}z (Unconstrained)}
\begin{itemize}
    \item \textbf{K{\i}s{\i}ts{\i}z:} $x$'e hi\c{c}bir s{\i}n{\i}r yok, istedi\u{g}ini se\c{c}
    \item \textbf{K{\i}s{\i}tl{\i}:} $x$ baz{\i} ko\c{s}ullar{\i} sa\u{g}lamal{\i} (\"orne\u{g}in $x \geq 0$, veya $x_1 + x_2 = 10$)
\end{itemize}

\subsubsection{Lineer vs Nonlineer}
\begin{itemize}
    \item \textbf{Lineer:} $f(x) = 3x_1 + 2x_2$ gibi (t\"um terimlerin derecesi 1)
    \item \textbf{Nonlineer:} $f(x) = x_1^2 + \sin(x_2)$ gibi (kareli, trigonometrik vs.)
\end{itemize}

\subsubsection{Konveks (Convex) vs Non-Convex}
\begin{itemize}
    \item \textbf{Konveks:} Fonksiyonun tek bir ``\c{c}ukuru'' var $\to$ minimum bulduysan, o \textbf{global} minimumd\"ur
    \item \textbf{Non-convex:} Birden fazla \c{c}ukur olabilir $\to$ buldu\u{g}un minimum sadece \textbf{lokal} olabilir
\end{itemize}

\subsubsection{Global vs Lokal \c{C}\"oz\"um}
\begin{itemize}
    \item \textbf{Global minimum:} T\"um $x$'ler i\c{c}inde en k\"u\c{c}\"uk $f(x)$ de\u{g}eri
    \item \textbf{Lokal minimum:} Yak{\i}n{\i}ndaki $x$'ler i\c{c}inde en k\"u\c{c}\"uk $f(x)$ de\u{g}eri
\end{itemize}

\textbf{\"Ornek:} $f(x) = x^4 - 2x^2 + 1$
\begin{itemize}
    \item $x = -1$ ve $x = 1$'de lokal (ve global) minimum var
    \item $x = 0$'da lokal maksimum var
\end{itemize}

\subsection{Optimization Tree (A\u{g}a\c{c} Yap{\i}s{\i})}

\begin{verbatim}
                    Optimization
                   /            \
            Continuous         Discrete
            /        \              \
    Unconstrained  Constrained   Integer Programming
         |          /    |    \
     (Line search, Linear  Network  Stochastic
      Trust region) Prog.  Prog.    Prog.
\end{verbatim}

Bu derste \"o\u{g}renece\u{g}imiz y\"ontemler:
\begin{enumerate}
    \item \textbf{K{\i}s{\i}ts{\i}z optimizasyon:} Line search, Trust region, Conjugate gradient, Quasi-Newton
    \item \textbf{Lineer programlama:} Simplex, Interior point
    \item \textbf{K{\i}s{\i}tl{\i} optimizasyon:} Quadratic programming, Penalty methods, Active set
\end{enumerate}

\newpage

%% ============================================================
%% SECTION 2
%% ============================================================
\section{Konveks Kavramlar (Affine, Cone, Convex)}

\subsection{Kombinasyonlar -- Neden \"Onemli?}

Optimizasyonda ``\c{c}\"oz\"um uzay{\i}n{\i}n \c{s}ekli'' \c{c}ok \"onemli. E\u{g}er \c{c}\"oz\"um uzay{\i}n konveks ise, i\c{s}imiz \c{c}ok kolayla\c{s}{\i}yor \c{c}\"unk\"u lokal minimum = global minimum olur.

\"Oncelikle \textbf{kombinasyon} kavram{\i}ndan ba\c{s}layal{\i}m.

\subsection{Lineer Kombinasyon}

Elimizde vekt\"orler var: $x_1, x_2, \ldots, x_p$ (bunlar $\mathbb{R}^n$'de noktalar)

\[
x = \lambda_1 x_1 + \lambda_2 x_2 + \cdots + \lambda_p x_p
\]

Katsay{\i}lar ($\lambda_i$) herhangi bir reel say{\i} olabilir: pozitif, negatif, s{\i}f{\i}r\ldots

\textbf{\"Ornek (2D'de):} $x_1 = (1, 0)$, $x_2 = (0, 1)$
\begin{itemize}
    \item $3x_1 + 2x_2 = (3, 2)$ $\checkmark$ lineer kombinasyon
    \item $-x_1 + 5x_2 = (-1, 5)$ $\checkmark$ lineer kombinasyon
    \item Sonu\c{c}: T\"um $\mathbb{R}^2$ d\"uzlemini kapsars{\i}n
\end{itemize}

\subsection{Affine Kombinasyon}

Lineer kombinasyonla ayn{\i}, AMA katsay{\i}lar{\i}n toplam{\i} 1 olmal{\i}:

\[
\sum_{i=1}^{p} \lambda_i = 1
\]

\textbf{\"Ornek:} $x_1 = (1, 0)$, $x_2 = (0, 1)$
\begin{itemize}
    \item $0.5 \cdot x_1 + 0.5 \cdot x_2 = (0.5, 0.5)$ $\checkmark$ ($0.5 + 0.5 = 1$)
    \item $2 \cdot x_1 + (-1) \cdot x_2 = (2, -1)$ $\checkmark$ ($2 + (-1) = 1$)
    \item $3 \cdot x_1 + 2 \cdot x_2 = (3, 2)$ $\times$ ($3 + 2 = 5 \neq 1$)
\end{itemize}

\textbf{Geometrik anlam{\i}:} \.{I}ki nokta aras{\i}ndaki \textbf{do\u{g}ru} (ve onun uzant{\i}s{\i}). Yani iki noktadan ge\c{c}en sonsuz do\u{g}ru.

\subsection{Konik (Conical) Kombinasyon}

Katsay{\i}lar \textbf{negatif olamaz:}

\[
\lambda_i \geq 0
\]

\textbf{\"Ornek:} $x_1 = (1, 0)$, $x_2 = (0, 1)$
\begin{itemize}
    \item $3 \cdot x_1 + 2 \cdot x_2 = (3, 2)$ $\checkmark$ (ikisi de $\geq 0$)
    \item $-1 \cdot x_1 + 5 \cdot x_2 = (-1, 5)$ $\times$ ($-1 < 0$)
\end{itemize}

\textbf{Geometrik anlam{\i}:} Orijinden ba\c{s}layan bir \textbf{koni} (u\c{c}) \c{s}eklinde b\"olge.

\subsection{Konveks (Convex) Kombinasyon}

\textbf{Hem} katsay{\i}lar{\i}n toplam{\i} 1, \textbf{hem} hepsi $\geq 0$:

\[
\sum_{i=1}^{p} \lambda_i = 1 \quad \text{ve} \quad \lambda_i \geq 0
\]

\textbf{\"Ornek:} $x_1 = (1, 0)$, $x_2 = (0, 1)$
\begin{itemize}
    \item $0.3 \cdot x_1 + 0.7 \cdot x_2 = (0.3, 0.7)$ $\checkmark$
    \item $0.5 \cdot x_1 + 0.5 \cdot x_2 = (0.5, 0.5)$ $\checkmark$ (tam orta nokta)
    \item $2 \cdot x_1 + (-1) \cdot x_2$ $\times$ ($-1 < 0$)
\end{itemize}

\textbf{Geometrik anlam{\i}:} \.{I}ki nokta aras{\i}ndaki \textbf{do\u{g}ru par\c{c}as{\i}}. \"U\c{c} nokta i\c{c}in u\c{c}u birle\c{s}tiren \textbf{\"u\c{c}gen}.

\subsection{\"Ozet Tablosu}

\begin{table}[h]
\centering
\begin{tabular}{@{}lll@{}}
\toprule
\textbf{Kombinasyon} & \textbf{Ko\c{s}ul} & \textbf{Geometrik Anlam} \\
\midrule
Lineer   & yok                                     & T\"um uzay \\
Affine   & $\sum \lambda_i = 1$                    & Noktalardan ge\c{c}en do\u{g}ru/d\"uzlem \\
Konik    & $\lambda_i \geq 0$                      & Orijinden koni \\
Konveks  & $\sum \lambda_i = 1$ ve $\lambda_i \geq 0$ & Noktalar{\i}n i\c{c}indeki b\"olge \\
\bottomrule
\end{tabular}
\end{table}

\subsection{Konveks Fonksiyon}

Bir fonksiyon $f$ \textbf{konveks} ise, herhangi iki noktas{\i} aras{\i}na \c{c}izdi\u{g}in do\u{g}ru, fonksiyonun \textbf{\"ust\"unde} kal{\i}r:

\[
f(\alpha x + (1-\alpha)y) \leq \alpha f(x) + (1-\alpha)f(y) \quad \forall \alpha \in [0,1]
\]

\textbf{Basit\c{c}e:} \c{C}ukur \c{s}eklinde (U harfi gibi). Herhangi iki noktay{\i} birle\c{s}tiren do\u{g}ru, e\u{g}rinin alt{\i}na d\"u\c{s}mez.

\textbf{\"Ornekler:}
\begin{itemize}
    \item $f(x) = x^2$ $\to$ konveks $\checkmark$ (parabol yukar{\i} bak{\i}yor)
    \item $f(x) = |x|$ $\to$ konveks $\checkmark$ (V \c{s}eklinde)
    \item $f(x) = e^x$ $\to$ konveks $\checkmark$
    \item $f(x) = -x^2$ $\to$ konveks de\u{g}il, \textbf{konkav} (tepe noktas{\i} var)
\end{itemize}

\textbf{Konkav fonksiyon:} $-f$ konveks ise $f$ konkavd{\i}r. (Tepe \c{s}eklinde, ters U)

\textbf{Neden \"onemli?} Konveks fonksiyonun minimumunu bulursan, o \textbf{kesinlikle global minimumd\"ur}. Ba\c{s}ka bir yerde daha k\"u\c{c}\"uk de\u{g}er yoktur!

\newpage

%% ============================================================
%% SECTION 3
%% ============================================================
\section{Kalk\"ul\"us Tekrar{\i} (Optimizasyon \.{I}\c{c}in Gereken)}

\subsection{Yakla\c{s}{\i}m H{\i}z{\i} (Rate of Convergence)}

Optimizasyon algoritmalar{\i} \textbf{iteratif} \c{c}al{\i}\c{s}{\i}r: bir ba\c{s}lang{\i}\c{c} noktas{\i} se\c{c}, ad{\i}m ad{\i}m \c{c}\"oz\"ume yakla\c{s}.
Peki ne kadar h{\i}zl{\i} yakla\c{s}{\i}yorsun? Bunu \"ol\c{c}en kavram \textbf{yakla\c{s}{\i}m h{\i}z{\i}}.

$\{x_k\}$ dizisi $x^*$'a yakla\c{s}{\i}yor olsun:

\subsubsection{Q-Linear (Do\u{g}rusal Yakla\c{s}{\i}m)}

\[
\frac{\|x_{k+1} - x^*\|}{\|x_k - x^*\|} \leq r \quad \text{, } r \in (0,1)
\]

Her ad{\i}mda hatay{\i} \textbf{sabit bir oranla} azalt{\i}rs{\i}n.

\textbf{\"Ornek:} $r = 0.5$ ise, her ad{\i}mda hata yar{\i}ya iner.
\begin{itemize}
    \item Hata: $1 \to 0.5 \to 0.25 \to 0.125 \to \ldots$
\end{itemize}

\subsubsection{Q-Superlinear (Do\u{g}rusal-\"Ust\"u)}

\[
\lim_{k \to \infty} \frac{\|x_{k+1} - x^*\|}{\|x_k - x^*\|} = 0
\]

Oran giderek k\"u\c{c}\"ul\"uyor, yani her ad{\i}m bir \"oncekinden \textbf{oransal olarak daha iyi}.

\subsubsection{Q-Order $p$ ($p$. Dereceden)}

\[
\frac{\|x_{k+1} - x^*\|}{\|x_k - x^*\|^p} \leq M
\]

$p = 2$ ise \textbf{kuadratik yakla\c{s}{\i}m}: hata her ad{\i}mda \textbf{karesi} kadar azal{\i}r.
\begin{itemize}
    \item Hata: $0.1 \to 0.01 \to 0.0001 \to 0.00000001$
    \item Newton metodu bunu ba\c{s}ar{\i}r!
\end{itemize}

\subsection{Lipschitz S\"ureklili\u{g}i}

Bir fonksiyon $f$, $\mathcal{N}$ k\"umesinde \textbf{Lipschitz s\"urekli} ise:

\[
\|f(x) - f(y)\| \leq L \|x - y\| \quad \forall x, y \in \mathcal{N}
\]

\textbf{Ne demek?} Fonksiyonun de\u{g}i\c{s}im h{\i}z{\i} \textbf{s{\i}n{\i}rl{\i}}. Yani $f$ ``\c{c}ok h{\i}zl{\i}'' de\u{g}i\c{s}emez.
$L$ sabiti ne kadar k\"u\c{c}\"ukse, fonksiyon o kadar ``yava\c{s}'' de\u{g}i\c{s}ir.

\textbf{\"Ornek:} $f(x) = 2x$
\begin{itemize}
    \item $|f(x) - f(y)| = |2x - 2y| = 2|x - y|$
    \item $L = 2$ ile Lipschitz s\"urekli $\checkmark$
\end{itemize}

\textbf{\"Ornek:} $f(x) = x^2$ (t\"um $\mathbb{R}$'de)
\begin{itemize}
    \item $|x^2 - y^2| = |x+y||x-y|$
    \item $|x+y|$ s{\i}n{\i}rl{\i} de\u{g}ilse $L$ bulunamaz $\to$ T\"um $\mathbb{R}$'de Lipschitz de\u{g}il
    \item Ama $[-5, 5]$ gibi s{\i}n{\i}rl{\i} bir aral{\i}kta Lipschitz ($L = 10$)
\end{itemize}

\textbf{\"Onemli \"ozellik:} $f$ ve $g$ Lipschitz ise $\to$ $f+g$ ve $f \cdot g$ de Lipschitz.

\subsection{K{\i}smi T\"urev ve Gradient}

\subsubsection{Tek De\u{g}i\c{s}kenli T\"urev}
$f(x) = x^3 \;\Rightarrow\; f'(x) = 3x^2$

\subsubsection{\c{C}ok De\u{g}i\c{s}kenli: K{\i}smi T\"urev (Partial Derivative)}

$f(x_1, x_2)$ gibi birden fazla de\u{g}i\c{s}kenli fonksiyonlarda, \textbf{bir de\u{g}i\c{s}kene g\"ore} t\"urev al{\i}rken di\u{g}er de\u{g}i\c{s}kenleri sabit tutars{\i}n.

\[
\frac{\partial f}{\partial x_i} = \lim_{h \to 0} \frac{f(x + h e_i) - f(x)}{h}
\]

\textbf{\"Ornek:} $f(x_1, x_2) = x_1^2 + 3x_1 x_2 + x_2^2$

\[
\frac{\partial f}{\partial x_1} = 2x_1 + 3x_2 \quad \text{($x_2$'yi sabit tutuyoruz)}
\]

\[
\frac{\partial f}{\partial x_2} = 3x_1 + 2x_2 \quad \text{($x_1$'i sabit tutuyoruz)}
\]

\subsubsection{Gradient (T\"urevlerin Vekt\"or\"u)}

T\"um k{\i}smi t\"urevleri bir vekt\"or i\c{c}inde toplarsak \textbf{gradient} elde ederiz:

\[
\nabla f(x) = \begin{pmatrix} \partial f / \partial x_1 \\ \partial f / \partial x_2 \\ \vdots \\ \partial f / \partial x_n \end{pmatrix}
\]

\textbf{Yukar{\i}daki \"ornek i\c{c}in:}

\[
\nabla f(x) = \begin{pmatrix} 2x_1 + 3x_2 \\ 3x_1 + 2x_2 \end{pmatrix}
\]

\textbf{Gradient ne i\c{s}e yarar?}
\begin{itemize}
    \item Gradient, fonksiyonun \textbf{en h{\i}zl{\i} artt{\i}\u{g}{\i} y\"on\"u} g\"osterir
    \item $-\nabla f$ (negatif gradient) ise \textbf{en h{\i}zl{\i} azald{\i}\u{g}{\i} y\"on} $\to$ minimuma gitmek i\c{c}in bu y\"onde y\"ur\"u!
\end{itemize}

\subsubsection{Y\"onl\"u T\"urev (Directional Derivative)}

Gradient herhangi bir $p$ y\"on\"undeki de\u{g}i\c{s}im h{\i}z{\i}n{\i} verir:

\[
D(f(x), p) = \nabla f(x)^T p
\]

Burada $\|p\| = 1$ (birim vekt\"or).

\subsection{Hessian Matrisi}

Hessian = \textbf{ikinci t\"urevlerin matrisi}. Tek de\u{g}i\c{s}kende $f''(x)$ ne ise, \c{c}ok de\u{g}i\c{s}kende Hessian odur.

\[
\nabla^2 f(x) = \begin{pmatrix} \dfrac{\partial^2 f}{\partial x_1^2} & \dfrac{\partial^2 f}{\partial x_1 \partial x_2} & \cdots \\[8pt] \dfrac{\partial^2 f}{\partial x_2 \partial x_1} & \dfrac{\partial^2 f}{\partial x_2^2} & \cdots \\[8pt] \vdots & \vdots & \ddots \end{pmatrix}
\]

\textbf{\"Ornek:} $f(x_1, x_2) = x_1^2 + 3x_1 x_2 + x_2^2$

\"Onceki k{\i}smi t\"urevlerden:
\begin{itemize}
    \item $\partial f / \partial x_1 = 2x_1 + 3x_2$
    \item $\partial f / \partial x_2 = 3x_1 + 2x_2$
\end{itemize}

\c{S}imdi bunlar{\i}n da t\"urevini alal{\i}m:

\[
\nabla^2 f = \begin{pmatrix} 2 & 3 \\ 3 & 2 \end{pmatrix}
\]

\textbf{\"Onemli:} $f$ iki kez t\"urevlenebilir ise Hessian \textbf{simetriktir} (sa\u{g} \"ust = sol alt).

\textbf{Hessian ne i\c{s}e yarar?}
\begin{itemize}
    \item Fonksiyonun \textbf{e\u{g}rili\u{g}ini} g\"osterir
    \item Hessian pozitif definit ise $\to$ o nokta \textbf{minimum} (\c{c}ukur)
    \item Hessian negatif definit ise $\to$ o nokta \textbf{maksimum} (tepe)
    \item Kar{\i}\c{s}{\i}k ise $\to$ \textbf{eyer noktas{\i}} (ne min ne max)
\end{itemize}

\subsection{Taylor Teoremi}

Tek de\u{g}i\c{s}kende biliyorsun:
\[
f(x+h) \approx f(x) + f'(x)h + \frac{1}{2}f''(x)h^2
\]

\c{C}ok de\u{g}i\c{s}kende \textbf{ayn{\i} mant{\i}k}, ama t\"urev yerine gradient, $f''$ yerine Hessian kullan{\i}l{\i}r:

\subsubsection{Ortalama De\u{g}er Teoremi (Mean Value Theorem)}
\[
f(x+p) = f(x) + \nabla f(x + \alpha p)^T p \quad \text{, } \alpha \in (0,1)
\]

\subsubsection{Taylor A\c{c}{\i}l{\i}m{\i} (2. derece)}
\[
f(x+p) = f(x) + \nabla f(x)^T p + \frac{1}{2} p^T \nabla^2 f(x + \alpha p) \, p
\]

\textbf{Neden \"onemli?} Optimizasyon algoritmalar{\i} fonksiyonu \textbf{Taylor a\c{c}{\i}l{\i}m{\i} ile yakla\c{s}t{\i}r{\i}p} minimum bulur. \"Orne\u{g}in Newton metodu tam olarak bu 2. derece yakla\c{s}{\i}m{\i} kullan{\i}r.

\subsection{Jacobian (Vekt\"or De\u{g}erli Fonksiyonlar \.{I}\c{c}in)}

\c{S}imdi $f$ tek bir say{\i} d\"ondermiyorsa, mesela $r: \mathbb{R}^n \to \mathbb{R}^m$ ($n$ giri\c{s}, $m$ \c{c}{\i}k{\i}\c{s}):

\[
r(x) = \begin{pmatrix} f_1(x_1, \ldots, x_n) \\ f_2(x_1, \ldots, x_n) \\ \vdots \\ f_m(x_1, \ldots, x_n) \end{pmatrix}
\]

\textbf{Jacobian} = her bir $f_i$'nin gradientini sat{\i}rlara yaz:

\[
J(x) = \begin{pmatrix} \nabla f_1(x)^T \\ \nabla f_2(x)^T \\ \vdots \\ \nabla f_m(x)^T \end{pmatrix}
\]

\textbf{\"Ornek:} $r(x_1, x_2) = (x_1^2 + x_2,\; x_1 x_2)$

\[
J = \begin{pmatrix} 2x_1 & 1 \\ x_2 & x_1 \end{pmatrix}
\]

\textbf{Gradient vs Jacobian:}
\begin{itemize}
    \item Gradient: skaler fonksiyon ($\mathbb{R}^n \to \mathbb{R}$) i\c{c}in, sonu\c{c} $n \times 1$ vekt\"or
    \item Jacobian: vekt\"or fonksiyon ($\mathbb{R}^n \to \mathbb{R}^m$) i\c{c}in, sonu\c{c} $m \times n$ matris
    \item Asl{\i}nda gradient, Jacobian'{\i}n \"ozel hali ($m=1$ durumu)
\end{itemize}

\newpage

%% ============================================================
%% SECTION 4
%% ============================================================
\section{Lineer Cebir Tekrar{\i}}

\subsection{Vekt\"or}

Bir \textbf{kolon vekt\"or} $x \in \mathbb{R}^n$:

\[
x = \begin{pmatrix} x_1 \\ x_2 \\ \vdots \\ x_n \end{pmatrix}
\]

\textbf{Transpose (devrik):} Kolon vekt\"or\"u sat{\i}r vekt\"ore \c{c}evirir:

\[
x^T = (x_1, x_2, \ldots, x_n)
\]

\subsubsection{\.{I}\c{c} \c{C}arp{\i}m (Inner Product / Dot Product)}

\[
x^T y = \sum_{i=1}^{n} x_i y_i = x_1 y_1 + x_2 y_2 + \cdots + x_n y_n
\]

\textbf{\"Ornek:} $x = (1, 2, 3)^T$, $y = (4, 5, 6)^T$
\begin{itemize}
    \item $x^T y = 1 \cdot 4 + 2 \cdot 5 + 3 \cdot 6 = 4 + 10 + 18 = 32$
\end{itemize}

\subsubsection{Vekt\"or Normlar{\i} (B\"uy\"ukl\"uk \"Ol\c{c}\"us\"u)}

Bir vekt\"or\"un ``uzunlu\u{g}unu'' \"ol\c{c}menin farkl{\i} yollar{\i}:

\begin{table}[h]
\centering
\begin{tabular}{@{}llll@{}}
\toprule
\textbf{Norm} & \textbf{Form\"ul} & \textbf{Anlam} & \textbf{\"Ornek: $x = (3, -4)$} \\
\midrule
1-norm & $\|x\|_1 = \sum |x_i|$ & Mutlak de\u{g}erlerin toplam{\i} & $|3| + |-4| = 7$ \\
2-norm & $\|x\|_2 = \sqrt{\sum x_i^2}$ & \"Oklid uzunlu\u{g}u & $\sqrt{9+16} = 5$ \\
$\infty$-norm & $\|x\|_\infty = \max |x_i|$ & En b\"uy\"uk eleman & $\max(3, 4) = 4$ \\
\bottomrule
\end{tabular}
\end{table}

\textbf{2-norm en \c{c}ok kullan{\i}lan.} ``Vekt\"or\"un uzunlu\u{g}u'' denince genelde budur.

\subsection{Matris}

$m \times n$ boyutlu matris $A$:

\[
A = \begin{pmatrix} A_{11} & A_{12} & \cdots & A_{1n} \\ A_{21} & A_{22} & \cdots & A_{2n} \\ \vdots & & & \vdots \\ A_{m1} & A_{m2} & \cdots & A_{mn} \end{pmatrix}
\]

\subsubsection{Transpose (Devrik)}
Sat{\i}rlar{\i} ve s\"utunlar{\i} yer de\u{g}i\c{s}tirir: $(A^T)_{ij} = A_{ji}$

\subsubsection{Simetrik Matris}
$A^T = A$ ise matris \textbf{simetriktir}. Yani $A_{ij} = A_{ji}$.

\textbf{\"Ornek:}
\[
A = \begin{pmatrix} 1 & 3 \\ 3 & 2 \end{pmatrix} \quad \text{simetriktir (sa\u{g} \"ust = sol alt)}
\]

\subsubsection{Matris Normu}

\[
\|A\|_p = \max_{\|x\|_p = 1} \|Ax\|_p
\]

``$A$'n{\i}n bir birim vekt\"or\"u ne kadar b\"uy\"utebilece\u{g}i''

\subsection{\"Ozde\u{g}er ve \"Ozvekt\"or (Eigenvalue \& Eigenvector)}

Bu kavram optimizasyonda \textbf{\c{c}ok} \"onemli!

\subsubsection{Tan{\i}m}
$A$ bir $n \times n$ matris. E\u{g}er s{\i}f{\i}r olmayan bir $x$ vekt\"or\"u i\c{c}in:

\[
Ax = \lambda x
\]

ise, $\lambda$ \textbf{\"ozde\u{g}er}, $x$ \textbf{\"ozvekt\"or}.

\textbf{Ne demek?} $A$ matrisi $x$ vekt\"or\"un\"u \c{c}arpt{\i}\u{g}{\i}nda, $x$'in \textbf{y\"on\"u de\u{g}i\c{s}miyor}, sadece $\lambda$ kat{\i} b\"uy\"uyor/k\"u\c{c}\"ul\"uyor.

\textbf{\"Ornek:}
\[
A = \begin{pmatrix} 2 & 1 \\ 1 & 2 \end{pmatrix}
\]

$x = (1, 1)^T$ deneyelim:
\[
Ax = \begin{pmatrix} 2 & 1 \\ 1 & 2 \end{pmatrix} \begin{pmatrix} 1 \\ 1 \end{pmatrix} = \begin{pmatrix} 3 \\ 3 \end{pmatrix} = 3 \begin{pmatrix} 1 \\ 1 \end{pmatrix}
\]

$\lambda = 3$, $x = (1,1)^T$ bir \"ozde\u{g}er-\"ozvekt\"or \c{c}ifti!

$x = (1, -1)^T$ deneyelim:
\[
Ax = \begin{pmatrix} 2 & 1 \\ 1 & 2 \end{pmatrix} \begin{pmatrix} 1 \\ -1 \end{pmatrix} = \begin{pmatrix} 1 \\ -1 \end{pmatrix} = 1 \begin{pmatrix} 1 \\ -1 \end{pmatrix}
\]

$\lambda = 1$, $x = (1,-1)^T$ bir di\u{g}er \c{c}ift!

\subsubsection{\"Ozde\u{g}erleri Bulma}
$\det(A - \lambda I) = 0$ denklemini \c{c}\"oz.

Yukar{\i}daki \"ornek i\c{c}in:
\begin{itemize}
    \item $\det\!\bigl(\begin{smallmatrix} 2-\lambda & 1 \\ 1 & 2-\lambda \end{smallmatrix}\bigr) = (2-\lambda)^2 - 1 = 0$
    \item $\lambda^2 - 4\lambda + 3 = 0 \;\Rightarrow\; \lambda = 3,\; \lambda = 1$ $\checkmark$
\end{itemize}

\subsubsection{Simetrik Pozitif Definit (SPD) Matris}

$A$ \textbf{simetrik pozitif definit} ise:
\begin{enumerate}
    \item $A^T = A$ (simetrik)
    \item T\"um \"ozde\u{g}erleri \textbf{pozitif} ($\lambda > 0$)
\end{enumerate}

\textbf{E\c{s}de\u{g}er tan{\i}m:} Her s{\i}f{\i}r olmayan $x$ i\c{c}in: $x^T A x > 0$

\textbf{Neden \"onemli?} $f(x)$'in Hessian'{\i} SPD ise $\to$ o nokta kesinlikle \textbf{minimum}!

\subsubsection{Eigen-Decomposition}

$A$'n{\i}n $n$ tane ba\u{g}{\i}ms{\i}z \"ozvekt\"or\"u varsa:
\[
A = X \Lambda X^{-1}
\]
\begin{itemize}
    \item $X$: \"ozvekt\"orler (s\"utunlarda)
    \item $\Lambda$: \"ozde\u{g}erler (k\"o\c{s}egen matris)
\end{itemize}

\subsection{Spectral Decomposition (Simetrik Matrisler \.{I}\c{c}in)}

$A$ \textbf{simetrik} ise, \"ozvekt\"orler \textbf{ortogonal} se\c{c}ilir:

\[
A = Q \Lambda Q^T
\]

\begin{itemize}
    \item $Q$: ortogonal matris ($Q^T Q = QQ^T = I$) $\to$ \"ozvekt\"orler
    \item $\Lambda$: k\"o\c{s}egen matris $\to$ \"ozde\u{g}erler
\end{itemize}

\textbf{Ortogonal matris} ne demek? Sat{\i}rlar{\i} (ve s\"utunlar{\i}) birbirine dik ve birim uzunlukta.
Ters almak i\c{c}in sadece transpose yeterli: $Q^{-1} = Q^T$ (hesap i\c{c}in \c{c}ok kolay!)

\subsection{Tekil De\u{g}er Ayr{\i}\c{s}{\i}m{\i} (SVD)}

\textbf{Herhangi} bir $m \times n$ matris $A$ i\c{c}in (simetrik olmak zorunda de\u{g}il!):

\[
A = U \Sigma V^T
\]

\begin{itemize}
    \item $U$: $m \times m$ ortogonal matris (sol tekil vekt\"orler)
    \item $\Sigma$: $m \times n$ k\"o\c{s}egen matris (tekil de\u{g}erler, $\sigma_1 \geq \sigma_2 \geq \ldots \geq 0$)
    \item $V$: $n \times n$ ortogonal matris (sa\u{g} tekil vekt\"orler)
\end{itemize}

\textbf{Tekil de\u{g}erler =} $A^T A$'n{\i}n \"ozde\u{g}erlerinin karek\"okleri.

\textbf{Neden \"onemli?} SVD hemen hemen her yerde kullan{\i}l{\i}r: en k\"u\c{c}\"uk kareler, boyut indirgeme, pseudo-inverse\ldots

\subsection{LU Ayr{\i}\c{s}{\i}m{\i} (LU Decomposition)}

$Ax = b$ lineer sistemi \c{c}\"ozmek i\c{c}in:

\[
PA = LU
\]

\begin{itemize}
    \item $P$: perm\"utasyon matrisi (sat{\i}r de\u{g}i\c{s}tirme)
    \item $L$: alt \"u\c{c}gensel matris (lower triangular)
    \item $U$: \"ust \"u\c{c}gensel matris (upper triangular)
\end{itemize}

\textbf{\c{C}\"oz\"um ad{\i}m ad{\i}m:}
\begin{enumerate}
    \item $PA = LU$ ayr{\i}\c{s}{\i}m{\i}n{\i} yap
    \item $Ly = Pb$ \c{c}\"oz (ileri yerine koyma -- kolay!)
    \item $Ux = y$ \c{c}\"oz (geri yerine koyma -- kolay!)
\end{enumerate}

\textbf{Neden do\u{g}rudan $A^{-1}b$ hesaplam{\i}yoruz?} Ters matris hesaplamak \c{c}ok pahal{\i} ($O(n^3)$). LU ayr{\i}\c{s}{\i}m{\i} da $O(n^3)$ ama daha verimli ve say{\i}sal olarak daha kararl{\i}.

\subsection{Cholesky Ayr{\i}\c{s}{\i}m{\i}}

$A$ \textbf{simetrik pozitif definit (SPD)} ise, LU'nun \"ozel hali:

\[
P^T A P = L L^T
\]

LU'ya g\"ore \textbf{2 kat h{\i}zl{\i}} \c{c}\"unk\"u simetriyi kullan{\i}yor.

$A$ SPD de\u{g}ilse \textbf{LBL ayr{\i}\c{s}{\i}m{\i}} kullan{\i}l{\i}r:

\[
P^T A P = L B L^T
\]

$B$: $1 \times 1$ veya $2 \times 2$ bloklu k\"o\c{s}egen matris.

\subsection{Alt Uzaylar ve QR Ayr{\i}\c{s}{\i}m{\i}}

\subsubsection{Null Uzay{\i} (Null Space)}
\[
\text{Null}(A) = \{w \mid Aw = 0,\; w \neq 0\}
\]
$A$'n{\i}n ``yok etti\u{g}i'' vekt\"orler k\"umesi.

\subsubsection{G\"or\"unt\"u (Range)}
\[
\text{Range}(A) = \{w \mid w = Av,\; \forall v\}
\]
$A$'n{\i}n ``\"uretebilece\u{g}i'' vekt\"orler k\"umesi.

\subsubsection{Temel Teorem}
\[
\text{Null}(A) \oplus \text{Range}(A^T) = \mathbb{R}^n
\]
Yani $\mathbb{R}^n$ uzay{\i}, $\text{Null}(A)$ ve $\text{Range}(A^T)$'nin do\u{g}rudan toplam{\i}na e\c{s}ittir.

\subsubsection{QR Ayr{\i}\c{s}{\i}m{\i}}
\[
AP = QR
\]
\begin{itemize}
    \item $Q$: ortogonal matris
    \item $R$: \"ust \"u\c{c}gensel matris
    \item $P$: perm\"utasyon matrisi
\end{itemize}

En k\"u\c{c}\"uk kareler problemlerinde \c{c}ok kullan{\i}l{\i}r.

\subsection{Singularite ve Ko\c{s}ul Say{\i}s{\i}}

\subsubsection{Singular (Tekil) Matris}
$A$ matrisi \textbf{tekil} (tersi yok) ise:
\begin{itemize}
    \item En az bir \"ozde\u{g}eri 0
    \item En az bir tekil de\u{g}eri 0
    \item Null uzay{\i} bo\c{s} de\u{g}il
    \item $\det(A) = 0$
\end{itemize}

\subsubsection{Ko\c{s}ul Say{\i}s{\i} (Condition Number)}

\[
\kappa(A) = \|A\| \cdot \|A^{-1}\|
\]

\begin{itemize}
    \item $\kappa(A)$ b\"uy\"ukse $\to$ $A$ \textbf{k\"ot\"ukonumlanm{\i}\c{s}} (ill-conditioned)
    \item K\"u\c{c}\"uk hatalar bile sonucu \c{c}ok de\u{g}i\c{s}tirebilir
    \item $\kappa(A) = 1$ ideal (ortogonal matrisler)
    \item $\kappa(A) = \infty$ ise $A$ tekil
\end{itemize}

\subsection{Sherman-Morrison-Woodbury Form\"ul\"u}

$A$'n{\i}n tersini biliyorsan ve $A$'ya \textbf{k\"u\c{c}\"uk bir g\"uncelleme} yaparsan, yeni tersi \textbf{h{\i}zl{\i}ca} hesaplayabilirsin.

\subsubsection{Rank-1 G\"uncelleme}
$\hat{A} = A + uv^T$ ise:

\[
\hat{A}^{-1} = A^{-1} - \frac{A^{-1} u v^T A^{-1}}{1 + v^T A^{-1} u}
\]

\subsubsection{Genel Durum}
$\hat{A} = A + UV^T$ ($U, V$: $n \times p$ matris) ise:

\[
\hat{A}^{-1} = A^{-1} - A^{-1} U (I + V^T A^{-1} U)^{-1} V^T A^{-1}
\]

\textbf{Neden \"onemli?} Quasi-Newton optimizasyon y\"ontemleri her ad{\i}mda Hessian'a rank-1 veya rank-2 g\"uncelleme yapar. Bu form\"ul sayesinde her seferinde tersi s{\i}f{\i}rdan hesaplamaya gerek kalmaz!

\newpage

%% ============================================================
%% SECTION 5
%% ============================================================
\section{\c{C}\"oz\"ulm\"u\c{s} \"Ornekler}

\subsection{\"Ornek 1: Gradient Hesaplama}

\textbf{Soru:} $f(x_1, x_2) = 2x_1^2 + x_1 x_2 + x_2^2$ fonksiyonunun gradientini bul ve $(1, 2)$ noktas{\i}nda hesapla.

\textbf{\c{C}\"oz\"um:}

K{\i}smi t\"urevleri alal{\i}m:

\[
\frac{\partial f}{\partial x_1} = 4x_1 + x_2
\]

\[
\frac{\partial f}{\partial x_2} = x_1 + 2x_2
\]

Gradient:

\[
\nabla f(x) = \begin{pmatrix} 4x_1 + x_2 \\ x_1 + 2x_2 \end{pmatrix}
\]

$(1, 2)$ noktas{\i}nda:

\[
\nabla f(1, 2) = \begin{pmatrix} 4(1) + 2 \\ 1 + 2(2) \end{pmatrix} = \begin{pmatrix} 6 \\ 5 \end{pmatrix}
\]

\begin{yorum}
$(1, 2)$ noktas{\i}nda fonksiyon en h{\i}zl{\i} $(6, 5)$ y\"on\"unde artar. Minimum bulmak i\c{c}in $(-6, -5)$ y\"on\"unde y\"ur\"umeliyiz.
\end{yorum}

\subsection{\"Ornek 2: Hessian Hesaplama ve Pozitif Definitlik}

\textbf{Soru:} $f(x_1, x_2) = 3x_1^2 - 2x_1 x_2 + x_2^2$ fonksiyonunun Hessian{\i}n{\i} bul. Pozitif definit mi?

\textbf{\c{C}\"oz\"um:}

Gradient:
\[
\nabla f = \begin{pmatrix} 6x_1 - 2x_2 \\ -2x_1 + 2x_2 \end{pmatrix}
\]

Hessian (gradientin t\"urevini al):
\[
\nabla^2 f = \begin{pmatrix} 6 & -2 \\ -2 & 2 \end{pmatrix}
\]

Pozitif definitlik kontrol\"u -- \"ozde\u{g}erlere bakal{\i}m:

$\det(H - \lambda I) = 0$:
\begin{align*}
(6-\lambda)(2-\lambda) - (-2)(-2) &= 0 \\
\lambda^2 - 8\lambda + 12 - 4 &= 0 \\
\lambda^2 - 8\lambda + 8 &= 0 \\
\lambda &= \frac{8 \pm \sqrt{64-32}}{2} = \frac{8 \pm \sqrt{32}}{2} \\
\lambda_1 &\approx 6.83, \quad \lambda_2 \approx 1.17
\end{align*}

\textbf{\.{I}ki \"ozde\u{g}er de pozitif $\to$ Hessian pozitif definit $\to$ $f$ konveks bir fonksiyon!}

\subsection{\"Ornek 3: Taylor A\c{c}{\i}l{\i}m{\i} ile Yakla\c{s}{\i}m}

\textbf{Soru:} $f(x) = x_1^2 + x_2^2$ fonksiyonunu $x = (1, 1)$ noktas{\i} etraf{\i}nda 2. derece Taylor ile yakla\c{s}.

\textbf{\c{C}\"oz\"um:}

\[
f(x+p) \approx f(x) + \nabla f(x)^T p + \frac{1}{2} p^T \nabla^2 f(x)\, p
\]

$f(1, 1) = 1 + 1 = 2$

\[
\nabla f(1, 1) = \begin{pmatrix} 2 \\ 2 \end{pmatrix}
\]

\[
\nabla^2 f = \begin{pmatrix} 2 & 0 \\ 0 & 2 \end{pmatrix}
\]

$p = (0.1, 0.1)$ ad{\i}m{\i} atarsak:

\begin{align*}
f(1.1, 1.1) &\approx 2 + (2, 2) \begin{pmatrix} 0.1 \\ 0.1 \end{pmatrix} + \frac{1}{2} (0.1, 0.1) \begin{pmatrix} 2 & 0 \\ 0 & 2 \end{pmatrix} \begin{pmatrix} 0.1 \\ 0.1 \end{pmatrix} \\
&= 2 + 0.4 + \frac{1}{2}(0.04) = 2 + 0.4 + 0.02 = 2.42
\end{align*}

Ger\c{c}ek de\u{g}er: $f(1.1, 1.1) = 1.21 + 1.21 = 2.42$ $\checkmark$ (Bu fonksiyon kuadratik oldu\u{g}u i\c{c}in Taylor tam sonu\c{c} verir)

\subsection{\"Ornek 4: Konvekslik Kontrol\"u}

\textbf{Soru:} $f(x) = e^x$ fonksiyonunun konveks oldu\u{g}unu g\"oster.

\textbf{\c{C}\"oz\"um:} Tan{\i}mdan: $f(\alpha x + (1-\alpha)y) \leq \alpha f(x) + (1-\alpha)f(y)$ mi?

\[
e^{\alpha x + (1-\alpha)y} \leq \alpha e^x + (1-\alpha) e^y
\]

\.{I}kinci t\"urev testi: $f''(x) = e^x > 0$, $\forall x$ $\checkmark$

$f''(x) > 0$ her yerde $\to$ $f$ kesinlikle konveks.

\textbf{Kural:} Tek de\u{g}i\c{s}kenli fonksiyon i\c{c}in $f''(x) \geq 0$ ise konveks. \c{C}ok de\u{g}i\c{s}kenli i\c{c}in Hessian pozitif semi-definit ise konveks.

\subsection{\"Ornek 5: \"Ozde\u{g}er Bulma}

\textbf{Soru:} $A = \begin{pmatrix} 4 & 2 \\ 2 & 1 \end{pmatrix}$ matrisinin \"ozde\u{g}erlerini ve \"ozvekt\"orlerini bul.

\textbf{\c{C}\"oz\"um:}

$\det(A - \lambda I) = 0$:

\begin{align*}
(4-\lambda)(1-\lambda) - 4 &= 0 \\
4 - 5\lambda + \lambda^2 - 4 &= 0 \\
\lambda^2 - 5\lambda &= 0 \\
\lambda(\lambda - 5) &= 0
\end{align*}

$\lambda_1 = 0$, $\lambda_2 = 5$

$\lambda_1 = 0$ i\c{c}in \"ozvekt\"or: $(A - 0I)x = 0$:

\[
\begin{pmatrix} 4 & 2 \\ 2 & 1 \end{pmatrix} \begin{pmatrix} x_1 \\ x_2 \end{pmatrix} = 0
\]

$4x_1 + 2x_2 = 0 \;\Rightarrow\; x_2 = -2x_1 \;\Rightarrow\;$ $v_1 = (1, -2)^T$

$\lambda_2 = 5$ i\c{c}in: $(A - 5I)x = 0$:

\[
\begin{pmatrix} -1 & 2 \\ 2 & -4 \end{pmatrix} \begin{pmatrix} x_1 \\ x_2 \end{pmatrix} = 0
\]

$-x_1 + 2x_2 = 0 \;\Rightarrow\; x_1 = 2x_2 \;\Rightarrow\;$ $v_2 = (2, 1)^T$

\textbf{Not:} $\lambda_1 = 0$ oldu\u{g}u i\c{c}in $A$ tekil (singular). $\det(A) = 4 \cdot 1 - 2 \cdot 2 = 0$ $\checkmark$

\subsection{\"Ornek 6: LU Ayr{\i}\c{s}{\i}m{\i}}

\textbf{Soru:} $A = \begin{pmatrix} 2 & 1 \\ 4 & 3 \end{pmatrix}$ i\c{c}in LU ayr{\i}\c{s}{\i}m{\i}n{\i} bul ve $Ax = b$'yi \c{c}\"oz ($b = (3, 7)^T$).

\textbf{\c{C}\"oz\"um:}

\textbf{Ad{\i}m 1: LU Ayr{\i}\c{s}{\i}m{\i}} (pivotlama gerekmiyorsa $P = I$)

Gauss eliminasyonu yapal{\i}m:
\begin{itemize}
    \item Sat{\i}r 2'den $(4/2 = 2) \times$ Sat{\i}r 1'i \c{c}{\i}kar
\end{itemize}

\[
L = \begin{pmatrix} 1 & 0 \\ 2 & 1 \end{pmatrix}, \quad U = \begin{pmatrix} 2 & 1 \\ 0 & 1 \end{pmatrix}
\]

Kontrol: $LU = \begin{pmatrix} 2 & 1 \\ 4 & 3 \end{pmatrix} = A$ $\checkmark$

\textbf{Ad{\i}m 2: $Ly = b$ \c{c}\"oz} (ileri yerine koyma)

\[
\begin{pmatrix} 1 & 0 \\ 2 & 1 \end{pmatrix} \begin{pmatrix} y_1 \\ y_2 \end{pmatrix} = \begin{pmatrix} 3 \\ 7 \end{pmatrix}
\]

\begin{itemize}
    \item $y_1 = 3$
    \item $2(3) + y_2 = 7 \;\Rightarrow\; y_2 = 1$
\end{itemize}

\textbf{Ad{\i}m 3: $Ux = y$ \c{c}\"oz} (geri yerine koyma)

\[
\begin{pmatrix} 2 & 1 \\ 0 & 1 \end{pmatrix} \begin{pmatrix} x_1 \\ x_2 \end{pmatrix} = \begin{pmatrix} 3 \\ 1 \end{pmatrix}
\]

\begin{itemize}
    \item $x_2 = 1$
    \item $2x_1 + 1 = 3 \;\Rightarrow\; x_1 = 1$
\end{itemize}

\textbf{\c{C}\"oz\"um: $x = (1, 1)^T$}

\subsection{\"Ornek 7: Jacobian Hesaplama}

\textbf{Soru:} $r(x_1, x_2) = (x_1^2 + x_2^2,\; x_1 x_2,\; e^{x_1})$ fonksiyonunun Jacobian{\i}n{\i} bul.

\textbf{\c{C}\"oz\"um:}

$r: \mathbb{R}^2 \to \mathbb{R}^3$ (2 giri\c{s}, 3 \c{c}{\i}k{\i}\c{s}), yani Jacobian $3 \times 2$ matris.

\[
J(x) = \begin{pmatrix} \dfrac{\partial f_1}{\partial x_1} & \dfrac{\partial f_1}{\partial x_2} \\[8pt] \dfrac{\partial f_2}{\partial x_1} & \dfrac{\partial f_2}{\partial x_2} \\[8pt] \dfrac{\partial f_3}{\partial x_1} & \dfrac{\partial f_3}{\partial x_2} \end{pmatrix} = \begin{pmatrix} 2x_1 & 2x_2 \\ x_2 & x_1 \\ e^{x_1} & 0 \end{pmatrix}
\]

$(1, 1)$ noktas{\i}nda:
\[
J(1, 1) = \begin{pmatrix} 2 & 2 \\ 1 & 1 \\ e & 0 \end{pmatrix}
\]

\subsection{\"Ornek 8: Ko\c{s}ul Say{\i}s{\i}}

\textbf{Soru:} $A = \begin{pmatrix} 1 & 0 \\ 0 & 0.001 \end{pmatrix}$ matrisinin ko\c{s}ul say{\i}s{\i}n{\i} bul (2-norm ile).

\textbf{\c{C}\"oz\"um:}

$A$ k\"o\c{s}egen matris, \"ozde\u{g}erleri: $\lambda_1 = 1$, $\lambda_2 = 0.001$

2-norm i\c{c}in:
\begin{itemize}
    \item $\|A\|_2 = \max$ \"ozde\u{g}er $= 1$
    \item $\|A^{-1}\|_2 = \max(1/\lambda_i) = 1/0.001 = 1000$
\end{itemize}

\[
\kappa(A) = 1 \times 1000 = 1000
\]

\begin{yorum}
$\kappa(A) = 1000$ olduk\c{c}a b\"uy\"uk $\to$ bu matris \textbf{ill-conditioned}. $Ax = b$ \c{c}\"ozerken k\"u\c{c}\"uk hatalar 1000 kat b\"uy\"uyebilir.
\end{yorum}

\subsection{\"Ornek 9: Konveks Kombinasyon}

\textbf{Soru:} $A = (1, 3)$ ve $B = (5, 1)$ noktalar{\i}n{\i}n konveks kombinasyonlar{\i}n{\i} yaz. $(0.5, 0.5)$ ve $(0.3, 0.7)$ i\c{c}in hesapla.

\textbf{\c{C}\"oz\"um:}

Konveks kombinasyon: $P = \lambda A + (1-\lambda)B$, $\lambda \in [0,1]$

$\lambda = 0.5$:
\begin{itemize}
    \item $P = 0.5(1,3) + 0.5(5,1) = (0.5+2.5,\; 1.5+0.5) = \mathbf{(3, 2)}$ (orta nokta)
\end{itemize}

$\lambda = 0.3$:
\begin{itemize}
    \item $P = 0.3(1,3) + 0.7(5,1) = (0.3+3.5,\; 0.9+0.7) = \mathbf{(3.8, 1.6)}$ ($B$'ye daha yak{\i}n)
\end{itemize}

\textbf{Geometrik anlam:} $\lambda \in [0,1]$ de\u{g}i\c{s}tik\c{c}e, $A$ ile $B$ aras{\i}ndaki do\u{g}ru par\c{c}as{\i} \"uzerinde t\"um noktalar{\i} elde ederiz.

\subsection{\"Ornek 10: Minimumu Bul (Gradient = 0)}

\textbf{Soru:} $f(x_1, x_2) = x_1^2 + x_2^2 - 2x_1 - 4x_2 + 5$ fonksiyonunun minimumunu bul.

\textbf{\c{C}\"oz\"um:}

\textbf{Ad{\i}m 1: Gradient $= 0$ yap}

\[
\nabla f = \begin{pmatrix} 2x_1 - 2 \\ 2x_2 - 4 \end{pmatrix} = \begin{pmatrix} 0 \\ 0 \end{pmatrix}
\]

\begin{itemize}
    \item $2x_1 - 2 = 0 \;\Rightarrow\; x_1 = 1$
    \item $2x_2 - 4 = 0 \;\Rightarrow\; x_2 = 2$
\end{itemize}

Kritik nokta: $(1, 2)$

\textbf{Ad{\i}m 2: Hessian ile kontrol}

\[
\nabla^2 f = \begin{pmatrix} 2 & 0 \\ 0 & 2 \end{pmatrix}
\]

\"Ozde\u{g}erleri: $\lambda_1 = 2$, $\lambda_2 = 2$ (ikisi de pozitif)

$\Rightarrow$ Hessian \textbf{pozitif definit} $\to$ bu nokta kesinlikle \textbf{minimum}!

\[
f(1, 2) = 1 + 4 - 2 - 8 + 5 = 0
\]

Minimum de\u{g}er $\mathbf{0}$, $(1, 2)$ noktas{\i}nda.

\end{document}
